\section{A two-register compiler and the binary lower bound}
\label{comp:section}

The relation alphabet has $2^{h^2}$ symbols. We now replace it by a
single fixed binary alphabet while keeping the nondeterministic
source linear in $h$. A direct monoid argument transfers the
relation information; it avoids charging a deterministic target for
remembering a position inside a long encoded block.

\subsection{Four letters on two cyclic registers}

For this section index the $h$ relation points by $\mathbb Z/h\mathbb Z$,
using representatives $0,\ldots,h-1$. Relabeling has no effect on
the results of Section~\ref{alg:section}. Take disjoint state sets
\[
 U=\{U_0,\ldots,U_{h-1}\},\qquad
 V=\{V_0,\ldots,V_{h-1}\}.
\]
Define four relations on $U\cup V$:
\begin{align}
 a&=\{(U_i,U_{i+1}):i\in\mathbb Z/h\mathbb Z\}\cup\Id_V,
       \label{comp:a}\\
 b&=\Id_U\cup\{(V_j,V_{j+1}):j\in\mathbb Z/h\mathbb Z\},
       \label{comp:b}\\
 t&=\Id_{U\cup V}\cup\{(U_0,V_0)\},
       \label{comp:t}\\
 r&=\{(V_j,U_j):j\in\mathbb Z/h\mathbb Z\}.
       \label{comp:r}
\end{align}
The letter $r$ is undefined on every state of $U$. Each of $a,b,r$
is partial deterministic; the only extra branch is the edge
$U_0\to V_0$ of $t$.

For an $h$-point relation $R$, let $R(x,y)\in\bits$ be its matrix
entry. Products below are read from left to right in the displayed
increasing order. Define
\begin{equation}\label{comp:four-code}
 w(R)=
 \left[
  \prod_{i=0}^{h-1}
   \left(
    \left[\prod_{j=0}^{h-1} t^{\,R(-i,-j)}b\right]a
   \right)
 \right]r .
\end{equation}
Here $t^0$ means the empty word and $t^1$ means the one-letter word.
The negative indices are reduced modulo $h$.

\begin{theorem}[Exact two-register compilation]\label{comp:register}
For every $x,y\in\mathbb Z/h\mathbb Z$, there is exactly one
path from $U_x$ to $U_y$ labeled $w(R)$ if $(x,y)\in R$, and
there is no such path otherwise. Every path starting in $U$
and surviving the complete word ends in $U$.
Furthermore,
\begin{equation}\label{comp:length}
 |w(R)|=h^2+h+1+|R|\le2h^2+h+1.
\end{equation}
\end{theorem}

\begin{proof}
Before crossing to $V$, a path starting at $U_x$ is at $U_{x+i}$
during row $i$: each completed row contributes one $a$, and all
$b$ letters fix $U$. It can use the additional branch of $t$
only in the unique row $i=-x$. Within that row it can cross at
column $j$ exactly when the displayed exponent $R(-i,-j)$ is one.

After crossing it is at $V_0$. The $b$ immediately following that
$t$, and the remaining $b$ letters in the row, move it $h-j$
places, leaving it at $V_{-j}$. Each later complete row contains
$h$ copies of $b$, so those rows have no further net effect on
its $V$ index; their $a$ letters fix $V$. All later $t$ letters
also fix that state. The final $r$ sends it to $U_{-j}$.
Thus a crossing in column $j=-y$ exists exactly when $R(x,y)=1$.

The row and column of such a crossing are unique. At every other
$t$ letter the path must choose the identity edge. There is
therefore exactly one successful path for each entry equal to one.
A path which never crosses stays in $U$ and is killed by $r$.
No surviving path ends in $V$, since all edges of $r$ end in $U$.
Finally the word has $h^2$ copies of $b$, $h$ copies of $a$,
one $r$, and $|R|$ copies of $t$, proving \eqref{comp:length}.
\end{proof}

\begin{figure}[tb]
\centering
\begin{tikzpicture}[>=Stealth,node distance=15mm,
  state/.style={draw=accent,rounded corners,minimum width=35mm,
    minimum height=10mm,align=center}]
\node[state] (U) {$U$ register\\row memory};
\node[state,right=18mm of U] (V) {$V$ register\\column memory};
\draw[->] (U) edge[bend left=18] node[above] {$t:\ U_0\to V_0$} (V);
\draw[->] (V) edge[bend left=18] node[below] {$r:\ V_j\to U_j$} (U);
\draw[->] (U) edge[loop left] node[left] {$a$ cycles} (U);
\draw[->] (V) edge[loop right] node[right] {$b$ cycles} (V);
\end{tikzpicture}
\caption{The four fixed relations. Identity edges of $a,b,t$ are
specified in \eqref{comp:a}--\eqref{comp:t}. The final $r$ has no
edge out of $U$, forcing exactly one successful crossing per block.}
\label{comp:register-figure}
\end{figure}

\begin{corollary}[Concatenations and multiplicity]\label{comp:products}
On the boundary states $U$, the relation of
$w(R_1)\cdots w(R_\ell)$ is exactly the Boolean product
$R_1\cdots R_\ell$. More strongly, its matrix of numbers of paths
between $U$ states is the ordinary nonnegative-integer matrix product
of the zero-one matrices of the $R_i$.
\end{corollary}

\begin{proof}
Theorem~\ref{comp:register} identifies each single-block path-count
matrix on $U$ with the corresponding zero-one matrix. Every
surviving path is in $U$ at each block boundary, so concatenation
uses the usual sum over intermediate states. The empty
concatenation gives the identity matrix. Taking the support of
the integer product gives the Boolean product.
\end{proof}

The exact path-count statement is stronger than the language
equivalence needed below. It provides a useful independent test:
spurious duplicate branches cannot hide behind Boolean
reachability. It also makes clear that the full NFA need not be
unambiguous. Concatenations can have several intermediate-state
paths between the same endpoints.

\subsection{A complete binary language with \texorpdfstring{$6h+2$}{6h+2} states}
\label{comp:binary}

Use the fixed code
\begin{equation}\label{comp:bit-code}
 a\mapsto00,\qquad b\mapsto01,\qquad
 t\mapsto10,\qquad r\mapsto11,
\end{equation}
and let $\enc(R)=\code(w(R))$.
For each $q\in U\cup V$ take three decoder states
\[
 (q,\eps),\quad(q,0),\quad(q,1).
\]
From $(q,\eps)$ the first bit $d$ goes to $(q,d)$; from $(q,d)$
the second bit $e$ executes the four-letter relation corresponding
to $de$ and returns to a boundary state $(q',\eps)$.
All these transitions move right.

Add one initial state $q_{\mathrm{in}}$ and one accepting state
$q_{\mathrm{acc}}$. On $\Lmark$, $q_{\mathrm{in}}$ guesses any
$(U_i,\eps)$ and moves right. At $\Rmark$, each $(U_i,\eps)$
has a stay move to $q_{\mathrm{acc}}$. All other endmarker
transitions are undefined, and $q_{\mathrm{acc}}$ has no successors.
No other state is accepting.

This specifies a language $B_h$ on \emph{every} binary string.
An odd-length string ends in a pending-bit state and rejects.
An even-length string is decoded in consecutive two-bit groups;
it accepts precisely when some resulting four-letter path starts
and ends in $U$. There is no promise that the word consists of
canonical relation blocks and no unstated treatment of malformed
blocks. Some noncanonical even-length words may accept, according
to this explicit transition table. The empty binary word accepts.

\begin{proposition}[Binary state and length counts]\label{comp:binary-count}
The specified automaton is a $6h+2$-state 1NFA and, for every
finite sequence of $h$-point relations,
\begin{equation}\label{comp:canonical-membership}
 \enc(R_1)\cdots\enc(R_\ell)\in B_h
 \quad\Longleftrightarrow\quad R_1\cdots R_\ell\ne0.
\end{equation}
The canonical block length is
\[
 |\enc(R)|=2(h^2+h+1+|R|)\le4h^2+2h+2.
\]
\end{proposition}

\begin{proof}
There are $2h$ relation states with three decoder copies each,
plus the two explicitly added states. The decoder executes exactly
one four-letter transition per two input bits. On a canonical
concatenation, Corollary~\ref{comp:products} gives precisely the
relation product on the $U$ boundary states. The initial and final
rules test whether this relation has any pair. This includes
$\ell=0$, since $\Id_U$ is nonempty. The length formula follows
from Theorem~\ref{comp:register}.
\end{proof}

The order of the block length cannot be improved for an injective
binary encoding of all relations. There are $2^{h^2}$ distinct
relations, but only $2^{L+1}-1$ binary strings of length at most $L$.
Thus the largest block length of any such encoding is at least
$h^2$. In this sense the quadratic maximum length is optimal up
to a constant factor. This counting observation does not imply
optimality of the number of compiler states.

\subsection{Singleton contexts force the full relation divisor}
\label{comp:quotient-section}

Let $\eta:\bits^*\to M$ be a monoid homomorphism into any monoid
which recognizes $B_h$: assume there is a set $A\subseteq M$ such
that $w\in B_h$ exactly when $\eta(w)\in A$. This includes
recognition by ordinary-symbol matching diagrams, since the two
endmarker diagrams and the test ports define the set $A$.

Let $S\subseteq M$ be the submonoid generated by
$\{\eta(\enc(R)):R\in\Rel{[h]}\}$. It contains the empty product,
with the same identity as $M$. We only require a relation quotient
on $S$, not on all of $M$.

\begin{lemma}[Encoded context separation]\label{comp:quotient}
There is a well-defined unital surjective homomorphism
\[
 \phi:S\twoheadrightarrow\Rel{[h]},\qquad
 \eta(\enc(R_1)\cdots\enc(R_\ell))
       \longmapsto R_1\cdots R_\ell .
\]
\end{lemma}

\begin{proof}
Suppose two canonical concatenations have the same element of
$M$, but their relation products $R$ and $T$ differ. Choose
$(i,j)$ in their symmetric difference and put
$E_i=\{(i,i)\}$ and $E_j=\{(j,j)\}$. Then
\[
 E_i R E_j\ne0\quad\Longleftrightarrow\quad(i,j)\in R,
\]
and the corresponding equivalence holds for $T$.
Surround each canonical concatenation by $\enc(E_i)$ on the left
and $\enc(E_j)$ on the right. By
\eqref{comp:canonical-membership}, precisely one resulting binary
word belongs to $B_h$. Yet multiplicativity of $\eta$ makes their
images equal, contradicting recognition. This proves
well-definedness, including comparisons with an empty
concatenation.

Concatenating representatives proves multiplicativity. The empty
word is sent to $\Id_{[h]}$, so the map is unital. Every relation
is the image of its single canonical block, proving surjectivity.
\end{proof}

This is the usual singleton-context syntactic argument, specialized
to the encoded submonoid; compare
\cite{Pin1997,OpenAILiveness2026}. A decoder need not label every
arbitrary binary word by an $h$-point relation in a multiplicative
way. The narrower generated submonoid is both sufficient and
necessary for the argument as stated.

\subsection{Proof and quantitative comparison}
\label{comp:main-proof}

\begin{proof}[Proof of Theorem~\ref{intro:main}]
Take the binary language from Proposition~\ref{comp:binary-count}
and an equivalent $s$-state 2DFA. By
Theorem~\ref{mach:representation}, its ordinary-symbol matching
diagrams lie in $\Br{k}$ for some $k\le8s+2$ and recognize $B_h$
with fixed endmarker diagrams and fixed test ports.
Lemma~\ref{comp:quotient} supplies an identity-containing
submonoid with a unital surjection onto $\Rel{[h]}$.
Theorem~\ref{alg:divisor} therefore gives
$8s+2\ge k\ge F(h)$, proving \eqref{intro:h-bound}.

For $n\ge14$, put $h=\lfloor(n-2)/6\rfloor\ge2$. The source has
$6h+2\le n$ states and may be padded with unreachable,
nonaccepting states to reach exactly $n$. Integer floor
arithmetic gives
\[
 \left\lfloor\frac{\lfloor(n-2)/6\rfloor-2}{9}\right\rfloor
 =\left\lfloor\frac{n-14}{54}\right\rfloor.
\]
Solving \eqref{intro:h-bound} for $s$ proves
\eqref{intro:n-bound}. The exponential asymptotic follows by
absorbing the floor and fixed multiplicative factors into $O(1)$.
\end{proof}

The integer recurrence of Proposition~\ref{alg:finite-recurrence}
gives the useful finite refinement
\begin{equation}\label{comp:integer-bound}
 s\ge
 \max\left\{1,\left\lceil\frac{D(h)-1}{4}\right\rceil\right\}.
\end{equation}
The additional $1$ simply records that a recognizer has at least
one state. The supplied recurrence script computes $D(h)$
using exact integers; no floating-point rounding enters
\eqref{comp:integer-bound}.

\begin{table}[tb]
\centering
\small
\begin{tabular}{p{0.26\textwidth}p{0.31\textwidth}p{0.31\textwidth}}
\toprule
Quantity & Inspected earlier construction & This article\\
\midrule
Relation-divisor degree
& $2^{\lfloor(h-2)/31\rfloor}$
& $2(5/2)^{\lfloor(h-2)/9\rfloor}$\\[3pt]
2DFA matching degree
& $4(s+1)^2$ for arbitrary alphabets
& $2(a+2)s+2$ for $a$ letters\\[3pt]
Binary matching degree
& general bound $4(s+1)^2$
& $8s+2$\\[3pt]
Four-letter relation compiler
& $3h$ states in R816's convention
& $2h$ relation states\\[3pt]
Binary source count
& $9h$ states in R816's convention
& $6h+2$ including our marker states\\[3pt]
Binary conclusion
& constant-factor transfer applies to an imported liveness bound
& explicit bound \eqref{intro:n-bound}\\
\bottomrule
\end{tabular}
\caption{Comparison with the source matching argument
\cite{OpenAILiveness2026} and the earlier encoding record
\cite{TheoremDBR816}. Machine conventions explain the two added
marker states in the present source count. The general quadratic
port bound can be preferable when the alphabet is very large.}
\label{comp:comparison}
\end{table}

The leading exponential rate in the algebraic degree estimate
increases from $1/31$ to $\log_2(5/2)/9$ when measured in $h$.
The sparse matching representation then avoids taking a square
root of that degree lower bound for a fixed alphabet. These
are separate improvements; the binary encoding permits both
to contribute to a bound measured in the number of source
states. The lower bound concerns states, not running time,
number of head reversals, or the length of a transition-table
description.
